\documentclass[a4paper,twoside]{article}

\usepackage{morewrites}

\usepackage[svgnames,dvipsnames,x11names]{xcolor}
\usepackage{graphicx}
\usepackage[english]{babel}
\usepackage[colorlinks=true, linkcolor=NavyBlue, urlcolor=NavyBlue, citecolor=NavyBlue]{hyperref}
\usepackage[a4paper]{geometry}
\usepackage{ts-fonts}
\usepackage{amsmath}
\usepackage{float}
\usepackage{seqsplit}

\usepackage{lastpage}
\usepackage{refcount}
\makeatletter
\def\ps@plain{
  \let\@oddhead\@empty
  \let\@evenhead\@empty
  \def\@oddfoot{
    \hfil
    \ifnum\getpagerefnumber{LastPage}>1\relax
      \thepage
    \fi
    \hfil}
  \let\@evenfoot\@oddfoot
}
\makeatother

\title{Arguments}
\author{}\date{}

\begin{document}

\maketitle

\thispagestyle{plain}
\section{Attribute ownership \& consumers}\label{attribute-ownership-consumers}

\begin{itemize}
\item{} \textbf{Owner}: each attribute belongs to either the template (default) or a fragment (\texttt{fragment.toml} attributes set \texttt{owner = <fragment name>} implicitly). Conflicting owners raise a \texttt{TemplateError}.
\item{} \textbf{Emitter}: templates can surface derived attributes through \texttt{emit} in \texttt{manifest.toml} so renderers see both declared attributes and emitted helpers (for example, \texttt{callout\_style} or \texttt{code.engine}).
\item{} \textbf{Consumer}: any template/fragment/renderer code that reads the attribute. Consumers should only read attributes they own or that are explicitly emitted.
\item{} Precedence: CLI/front matter overrides \(\rightarrow\) attribute resolver (type coercion, normaliser, escape) \(\rightarrow\) emitted defaults \(\rightarrow\) render-time context.

\end{itemize}
When adding new attributes, pick a single owner (template or fragment) and document the sources that are allowed to override it.

\end{document}
